\chapter{Guía de Operación}
\label{ch:ops}
% ==============================================================

Este capítulo es la guía \textbf{definitiva de uso} de la aplicación tal y como queda tras la
reunión de equipo del 11 de septiembre de 2026 (transectos, entrenamiento de la línea húmeda,
edición manual, revisar/descartar imágenes) — pensada para quien va a \textbf{usar} la app
instalada, no para quien la desarrolla. La sección~\ref{sec:ops-dev} recoge, aparte, el flujo de
trabajo para quien sí toca código.

\section{Abrir la aplicación}

Para el uso normal no hace falta Python, Node.js ni ninguna terminal: se instala el ejecutable
de Windows (capítulo~\ref{ch:setup}, \S\ref{sec:build-installer}) y se abre como cualquier otro
programa.

\begin{enumerate}[leftmargin=1.5em]
  \item Ejecutar \texttt{LineaDeCosta-Setup.exe} y seguir el asistente de instalación.
  \item Abrir \textbf{``Línea de Costa''} desde el menú Inicio o el acceso directo del
        escritorio.
  \item La primera vez puede tardar unos segundos en aparecer la ventana: el propio ejecutable
        arranca primero el backend en segundo plano y luego abre la ventana nativa
        (\emph{pywebview}) apuntando a él — no hace falta ningún navegador aparte.
\end{enumerate}

\begin{warnbox}[No usar el acceso directo ``Descargar modelo de IA (SAM)''para el uso normal]
El instalador crea, además del acceso a la app, uno secundario en el menú Inicio para
\emph{forzar} la descarga del checkpoint de SAM ($\sim$2,4~GB) sin abrir la ventana — solo hace
falta si SAM no se instaló bien la primera vez o se quiere volver a descargar aposta. Abrirlo por
error solo muestra una consola que imprime una línea y se cierra sola; no es un fallo, pero
confunde si se pulsa pensando que es la app.
\end{warnbox}

\section{Recorrido por pantallas}

El menú lateral agrupa las pantallas en dos bloques: \textbf{Proceso} (el trabajo diario) y
\textbf{Sistema} (configuración y mantenimiento, se toca con menos frecuencia).

\subsection{Dashboard}

Vista general: estado de calibración de las 6 cámaras, gráfica de evolución histórica del área
seca y accesos rápidos a cada cámara.

\subsection{Carga de imágenes}

Pantalla para gestionar qué imágenes tiene cada cámara y lanzar su análisis.

\begin{itemize}[leftmargin=1.5em]
  \item \textbf{Subir imágenes}: arrastrar archivos a la zona punteada, o hacer clic para elegir
        con el explorador.
  \item \textbf{Filtro de fechas}: acota la cuadrícula de miniaturas a un rango \emph{desde/hasta}.
  \item \textbf{Selección}: clic sobre una miniatura para marcarla/desmarcarla; \textit{Todas} /
        \textit{Ninguna} / \textit{Solo H:00h} (la hora con más fotos del filtro actual) la
        rellenan de golpe.
  \item \textbf{Previsualizar antes de procesar} (nuevo, 2026-09-14): el icono de lupa que
        aparece al pasar el ratón sobre una miniatura abre la foto real a tamaño completo —
        la miniatura es demasiado pequeña para juzgar si una imagen sirve (niebla, cámara
        movida, algo tapando la zona). Desde ese mismo visor se puede \textbf{Descartar}
        (elimina el archivo) o \textbf{Cerrar} sin tocar nada.
  \item \textbf{Punto verde} en una miniatura: esa imagen ya tiene un resultado aceptado en el
        histórico.
  \item \textbf{Botón de la papelera} (esquina inferior derecha de la miniatura, aparece al
        pasar el ratón): borra la imagen directamente, con confirmación.
  \item \textbf{Botones de proceso}, de izquierda a derecha:
    \begin{itemize}
      \item \textit{Reanalizar todas (N)} — repite el análisis de TODAS las imágenes del filtro
            actual, incluidas las ya analizadas. Hace falta tras una mejora de la segmentación,
            para refrescar resultados antiguos.
      \item \textit{Analizar pendientes (N)} — solo las que nunca se han analizado.
      \item \textit{Procesar selección (N)} — solo las marcadas a mano.
    \end{itemize}
        Las tres procesan las imágenes \textbf{una a una} (SAM no admite lote) y al terminar
        llevan a \textit{Resultados} con la última imagen procesada.
\end{itemize}

\subsection{Resultados}

Pantalla principal de análisis: seleccionar cámara e imagen, pulsar \textit{Analizar imagen} y
revisar la línea de costa detectada y sus métricas.

\begin{itemize}[leftmargin=1.5em]
  \item \textbf{Flechas de navegación} (junto al título): pasan a la imagen anterior/siguiente de
        la lista. Si esa imagen \textbf{ya tiene} un resultado aceptado en el histórico, se
        muestra al instante, sin volver a ejecutar SAM (\S\ref{sec:cache-por-imagen}); si nunca
        se analizó, se lanza un análisis real, igual que con el botón de siempre.
  \item \textbf{Analizar todas (N)}: analiza de golpe todas las imágenes de la cámara
        seleccionada, sin salir de esta pantalla — para ver la evolución completa sin ir a
        Carga de imágenes.
  \item \textbf{Panel de métricas}, de arriba a abajo:
    \begin{enumerate}
      \item \textbf{Confianza} — barra verde/ámbar/rojo.
      \item \textbf{Anchura de playa (transectos)} — un valor en metros por cada transecto
            marcado para esta cámara (\S\ref{sec:transectos}); \textbf{indicador principal},
            reemplaza al área en fiabilidad. No aparece si la cámara no tiene ningún transecto
            marcado todavía.
      \item \textbf{Área seca (m\textsuperscript{2})} — se mantiene visible como referencia, no
            se ha eliminado nada.
      \item \textbf{Puntos UTM} y \textbf{Timestamp}.
    \end{enumerate}
  \item \textbf{Editar línea manualmente} (nuevo, 2026-09-14, botón junto al título, solo
        visible con un resultado aceptado): corrige a mano la línea que dio SAM cuando se
        equivoca — ver \S\ref{sec:edicion-manual}.
  \item \textbf{Exportar GeoJSON}: descarga el histórico completo de la cámara en un
        \texttt{.geojson} (EPSG:25830), listo para abrir en QGIS.
\end{itemize}

\subsubsection{Caché por imagen: por qué las flechas no vuelven a lanzar SAM}
\label{sec:cache-por-imagen}

Cada análisis se guarda con su nombre de fichero, tanto la imagen con la línea dibujada
(\texttt{GET /api/cameras/\{id\}/analysis-result?file=...}) como el resultado numérico
(\texttt{GET /api/cameras/\{id\}/images/\{filename\}/cached-result}, que lee directamente el
histórico). Navegar con las flechas a una imagen ya analizada consulta ese resultado guardado en
vez de relanzar SAM — instantáneo, y no satura el log con análisis repetidos de fotos que no han
cambiado.

\subsubsection{Edición manual de la línea detectada}
\label{sec:edicion-manual}

\begin{infobox}[Pedido en la reunión del 11/09/2026]
``Permite editar manualmente las segmentaciones automáticas y conserva las correcciones como
datos de entrenamiento.''
\end{infobox}

Al pulsar \textit{Editar línea manualmente}, la imagen mostrada cambia de la visualización
(con la línea ya ``quemada'' en el JPG) a la \textbf{foto original}, con la línea detectada por
SAM superpuesta como puntos de control editables —\texttt{analyze\_roi()} simplifica la línea
píxel a píxel (que sale dentada) con \texttt{cv2.approxPolyDP} a un puñado de puntos manejables,
no cientos.

\begin{itemize}[leftmargin=1.5em]
  \item \textbf{Clic} sobre la imagen: añade un punto nuevo.
  \item \textbf{Arrastrar} un punto existente: lo mueve.
  \item \textbf{Doble clic} sobre un punto: lo borra.
  \item \textit{Guardar corrección}: envía los puntos corregidos al \textbf{mismo almacén} que
        usa la sección de Entrenamiento (\S\ref{sec:entrenamiento}) — una corrección hecha aquí y
        una polilínea marcada a mano en Entrenamiento acaban en el mismo conjunto de datos, no en
        dos sitios separados.
  \item \textit{Cancelar}: descarta los cambios sin guardar nada.
\end{itemize}

\subsection{Mapa GeoJSON}

Visualiza sobre un mapa real las líneas de costa detectadas de todas las cámaras, en EPSG:25830.

\subsection{Calibración}

Asistente de 7 pasos para establecer la homografía píxel~$\to$~UTM de cada cámara — detallado en
el capítulo~\ref{ch:calibration}. No hace falta repetirlo salvo que la cámara se desplace de
forma notable (viento fuerte, limpieza brusca) o se añada una cámara nueva.

\subsection{Cámaras}

Configuración de hardware por cámara: lente, ROI (región de interés) y ubicación.

\subsection{Modo automático}

Automatiza de punta a punta, para un rango de fechas, lo que \textit{Carga de imágenes} y
\textit{Resultados} hacen a mano: descarga de la API de Obscape, alineación y análisis —
capítulo~\ref{ch:automode}. Requiere una cámara ya calibrada.

\subsection{Configuración}

\begin{itemize}[leftmargin=1.5em]
  \item \textbf{Credenciales de Obscape}: usuario y clave de API para descargar imágenes — se
        guardan solo en este ordenador.
  \item \textbf{Ubicación de los datos}: rutas reales en disco de imágenes y calibración de este
        ordenador, con botón \textit{Copiar}.
  \item \textbf{Entrenamiento línea húmeda (temporal)} (nuevo, 2026-09-14): ruta donde se guardan
        las polilíneas y las imágenes adicionales de la sección de Entrenamiento
        (\S\ref{sec:entrenamiento}) — informativa, para saber dónde ir a buscar los ficheros o
        compartir la carpeta con el equipo (p.\,ej. por Google Drive, como se comentó en la
        reunión).
  \item \textbf{Borrar imágenes y resultados}: acción destructiva e irreversible (requiere
        escribir \texttt{BORRAR} para confirmar) — elimina imágenes, varillas marcadas e
        histórico de las 6 cámaras. \textbf{La calibración no se toca.}
\end{itemize}

\subsection{Transectos --- anchura de playa}
\label{sec:transectos}

\begin{infobox}[Pedido en la reunión del 11/09/2026]
``Olvidaros de la cuestión de los cálculos de áreas [...] como indicador interesante sí que
podría ser un transecto [...] la longitud que hay desde ese punto hasta la línea que habéis
segmentado. Ese indicador nos da la anchura de playa en el tramo concreto.''
\end{infobox}

Sustituye al área en m\textsuperscript{2} como indicador principal de evolución de la playa: un
punto de referencia \textbf{fijo} (p.\,ej. la valla de la duna) + la distancia hasta la línea de
costa detectada en cada análisis.

\begin{enumerate}[leftmargin=1.5em]
  \item Ir a \textbf{Transectos} en el menú lateral y elegir la cámara.
  \item Elegir una imagen de referencia (por defecto, la más reciente).
  \item \textbf{Clic} sobre la imagen para añadir un transecto — aparece en la lista de la
        izquierda con una etiqueta editable (p.\,ej. ``Valla duna norte'').
  \item \textit{Guardar cambios}.
\end{enumerate}

A partir de ahí, \textbf{cada análisis} de esa cámara (en \textit{Resultados}, \textit{Analizar
todas} o \textit{Modo automático}) calcula automáticamente la distancia de cada transecto hasta
la línea detectada y la muestra en el panel de métricas — no hace falta volver a marcarlo.

\begin{warnbox}[Cómo se calcula la distancia]
El backend recalcula la posición UTM del punto marcado con la homografía \textbf{vigente} en
cada análisis (no la guarda fija), así que un transecto marcado hoy sigue siendo válido aunque la
cámara se recalibre más adelante. La distancia es la del punto al \textbf{segmento más cercano}
de la polilínea de costa — en la práctica, para un punto de referencia fijo mirando hacia el mar,
equivale a la distancia perpendicular que se mediría a mano con una cinta métrica, sin necesidad
de que el usuario marque también un ángulo (se discutió en la reunión y se optó por esta
simplificación en vez de pedir un segundo clic de ángulo).
\end{warnbox}

Un mismo transecto puede reutilizarse en todas las imágenes futuras de esa cámara — la cámara es
fija (o casi, ver debate sobre viento/limpieza en el capítulo~\ref{ch:calibration}), así que el
mismo píxel de referencia vale para cualquier imagen suya.

\subsection{Entrenamiento línea húmeda --- sección temporal}
\label{sec:entrenamiento}

\begin{warnbox}[Sección temporal]
Pensada para que el equipo de campo marque a mano dónde está la línea húmeda real en imágenes
reales, y así entrenar/ajustar la segmentación automática más adelante
(\S\ref{sec:linea-humeda-pendiente}). No afecta a los resultados de \textit{Resultados}.
\end{warnbox}

\begin{enumerate}[leftmargin=1.5em]
  \item Ir a \textbf{Entrenamiento (temp.)} en el menú lateral y elegir la cámara.
  \item Elegir una imagen de la lista desplegable (o subir una propia, ver más abajo).
  \item \textbf{Clic} sobre la imagen para ir añadiendo los puntos de la polilínea, en orden,
        siguiendo la línea húmeda real.
  \item \textit{Guardar polilínea}.
\end{enumerate}

\subsubsection{Zoom, lupa y ajuste fino}

Mismo patrón que la Marcación de varillas de Calibración (capítulo~\ref{ch:calibration}):

\begin{itemize}[leftmargin=1.5em]
  \item \textbf{Completa / Doble zoom}: alterna entre la imagen ajustada al panel y el doble de
        tamaño (con scroll).
  \item \textbf{Zoom de la lupa} (×2 a ×15): recorte ampliado que sigue al cursor, para acertar
        el píxel exacto en fotos de varios miles de píxeles de ancho.
  \item \textbf{Clic sobre un punto ya puesto}: lo selecciona (se pone verde y más grande) sin
        moverlo.
  \item \textbf{Flechas del teclado}, con un punto seleccionado: lo desplazan 1~px (10~px con
        Mayús) — para el ajuste fino que el ratón no permite con precisión.
\end{itemize}

\subsubsection{Imágenes adicionales}

No hace falta que las imágenes a marcar sean las ya cargadas de la cámara: el botón \textit{Subir
imágenes adicionales} permite traer cualquier foto suelta solo para entrenamiento — se guardan en
una carpeta separada (\texttt{training\_images/}) y aparecen en la lista marcadas como
\texttt{[adicional]}. Se pueden eliminar de forma permanente con \textit{Eliminar esta imagen
adicional} (a diferencia de las imágenes de cámara, que solo se pueden \emph{descartar}, ver
abajo — borrarlas de verdad es cosa de \textit{Carga de imágenes}).

\subsubsection{Revisar y descartar antes de anotar}

\begin{infobox}[Pedido en la reunión del 11/09/2026]
``Previsualizar, revisar y descartar imágenes antes del procesamiento [...] igual que se
descartan los puntos de control.''
\end{infobox}

El botón \textit{Descartar esta imagen} saca la imagen de la cola de anotación por defecto — es
un descarte \textbf{blando y reversible}: no borra el archivo, igual que descartar una varilla de
calibración no borra el punto marcado, solo lo excluye del cálculo. El checkbox \textit{Mostrar
descartadas (N)} las vuelve a mostrar, con la opción de \textit{Reincluir}.

\subsubsection{Exportar}

\textit{Exportar JSON (cámara)} descarga todas las polilíneas guardadas de la cámara actual en un
único fichero — pensado para compartir con el equipo (p.\,ej. por Google Drive) y como base para
entrenar/ajustar la segmentación más adelante.

\section{Flujo de trabajo recomendado (día a día)}

\begin{enumerate}[leftmargin=1.5em, label=\textbf{\arabic*.}]
  \item \textit{Modo automático} (o \textit{Carga de imágenes} + \textit{Analizar pendientes})
        para procesar las imágenes nuevas de cada cámara ya calibrada.
  \item Revisar en \textit{Resultados} los rechazos y las confianzas bajas; si la línea
        detectada se equivoca de forma puntual, corregirla con \textit{Editar línea
        manualmente} en vez de descartar el análisis entero.
  \item Consultar la \textbf{anchura de playa por transectos} como indicador principal de
        evolución (una vez marcados los transectos de cada cámara, es automático a partir de
        ahí).
  \item Cuando el equipo de campo tenga tiempo: marcar polilíneas de línea húmeda en
        \textit{Entrenamiento} sobre un puñado de imágenes por cámara (objetivo de la reunión:
        $\sim$20), revisando/descartando antes las que sean ambiguas.
\end{enumerate}

\subsection{Lo que queda pendiente, a propósito}
\label{sec:linea-humeda-pendiente}

Cambiar el objetivo real de la segmentación de ``toda la arena seca'' a la \textbf{línea húmeda}
es el único punto de la reunión del 11/09/2026 que se deja explícitamente para más adelante: no
se puede programar a ciegas, depende de tener ya un conjunto real de imágenes anotadas
(\S\ref{sec:entrenamiento}) con el que decidir la referencia exacta antes de tocar el algoritmo
de segmentación.

\section{Para desarrolladores}
\label{sec:ops-dev}

El resto de esta sección recoge el flujo de trabajo para quien sí toca código — no hace falta
para el uso normal de la app instalada (secciones anteriores).

\subsection{Entorno de desarrollo}

\begin{codebox}[Arrancar backend + frontend en modo desarrollo]
\begin{minted}{bash}
# Windows (PowerShell) — activa venv/ y levanta backend+frontend con un
# workspace de desarrollo vacío en workspace/dev/ (no toca datos reales)
.\start_dev.ps1

# Linux/macOS (o si se prefiere Conda: ver environment.yml)
source venv/bin/activate   # o: conda activate cv-lit
./scripts/start_dev.sh

# Backend:  http://localhost:8000 (FastAPI + docs en /docs)
# Frontend: http://localhost:5173 (Vue 3 + Vite)
\end{minted}
\end{codebox}

\begin{codebox}[Workspace de demo — datos de muestra ya calibrados]
\begin{minted}{bash}
.\start_demo.ps1
# Siembra workspace/demo/ desde cero cada vez (imágenes de muestra +
# varillas del GCP.csv ya importadas y calibradas). No toca los datos reales.
\end{minted}
\end{codebox}

\subsection{Descarga de imágenes (CLI)}

\begin{codebox}[]
\begin{minted}{bash}
# Descarga estándar (últimas 2 semanas, 12:00h)
python3 acces_api/scheduled_download.py

# Diagnóstico de estado de cámaras
python3 scripts/verify_setup.py
\end{minted}
\end{codebox}

\subsection{Calibración}

\begin{warnbox}[Ya no hay herramienta de calibración por línea de comandos]
Las herramientas CLI que existían anteriormente (\texttt{calibration\_tool.py},
\texttt{visualizar\_calibracion.py}, \texttt{recalibrate.py}) se eliminaron durante el
endurecimiento del proyecto para producción: quedaban duplicadas con el asistente web y sin
mantenimiento. La calibración se hace exclusivamente desde la interfaz web, pestaña
\textit{Calibración} — asistente de 7 pasos, detallado en el capítulo~\ref{ch:calibration}.
Verificar el motor matemático de la homografía se hace ahora con
\texttt{python -m pytest tests/test\_homography\_math.py tests/test\_geometry\_warning.py}.
\end{warnbox}

\subsection{Alineación masiva (CLI)}

\begin{codebox}[]
\begin{minted}{bash}
# Alinear todas las imágenes de una carpeta respecto a la primera
python3 homography/align_images.py \
  --input  /ruta/imagenes/ \
  --output /ruta/salida/ \
  --cam 1                    # aplica máscaras SIFT de CAM_1

# Especificar imagen de referencia manualmente
python3 homography/align_images.py \
  --input  /ruta/imagenes/ \
  --output /ruta/salida/ \
  --ref    imagen_base.jpg \
  --cam 2
\end{minted}
\end{codebox}

Salidas:
\begin{itemize}[leftmargin=1.5em]
  \item \texttt{/output/aligned/} — imágenes alineadas
  \item \texttt{/output/diagnostics/} — visualizaciones de matches
  \item \texttt{/output/homographies.csv} — métricas y matrices $H$ por imagen
\end{itemize}

\subsection{Pipeline de extracción (CLI)}

\begin{codebox}[]
\begin{minted}{bash}
# Prototipo completo: ROI → Segmentación SAM → Línea de costa
python3 proces_images/test_mes3_pipeline.py --cam 1
\end{minted}
\end{codebox}

\subsection{Suite de tests}

\begin{codebox}[]
\begin{minted}{bash}
# Todo menos los tests de integración (necesitan la API real de Obscape)
python -m pytest tests/ --ignore=tests/test_api_integration.py -q
\end{minted}
\end{codebox}

% ==============================================================
